\SourceLeaf{352}
\begin{center}\large\bfseries [PROJECTION DES WÜRFELS.]\end{center}

\Needspace{8\baselineskip}\begin{center}[1.]\end{center}

Sind die complexen Werthe der orthographischen Projectionen von drei gleich langen und unter einander senkrechten Geraden
\[
a,\ b,\ c,
\]
so ist
\[
aa+bb+cc=0,
\]
z.\,B.
\[
a=1+8i,\qquad b=8+i,\qquad c=4-4i.
\]
Allgemein kann man setzen
\[
a=pp+qq,\qquad b=2ipq,\qquad c=ipp-iqq.
\]

\Needspace{9\baselineskip}\begin{center}[2.]\end{center}

Die zierlichste Form der Auflösung der allgemeinen Gleichung
\[
pp+qq+rr=0
\]
ist, $a$ und $b$ beliebige complexe Zahlen bedeutend,
\[
p=(a-b)(b-aii)
\]
\[
q=(b-bi)(aii-ai)
\]
\[
r=(bi-a)(ai-b).
\]

\Needspace{9\baselineskip}\begin{center}\bfseries BEMERKUNGEN.\end{center}

Ein Theil dieser beiden Notizen, die sich in einem Handbuche befinden und aus dem Jahre 1831 stammen, ist bereits in Bd.\,II, S.\,309 abgedruckt worden.

\indent Vom December 1839 bis Ostern 1840 hat \textsc{Gauss} eine Vorlesung über die \emph{Theorie der imaginären Grössen} gehalten, die E.\,\textsc{Heine} gehört und ausgearbeitet hat, wie bereits S.\,330 angeführt. Dieser Ausarbeitung sind die folgenden auf die Projection des Würfels bezüglichen Ausführungen entnommen.
\par\noindent\hfill\textsc{Stäckel}.\par
